\section{The finite-volume and annulus estimates use one response}
\label{sec:ex-annulus-same-data}

The finite-volume recording and the annulus bound must be attached to the
same analytic response. This is a short compatibility step, but it prevents
an invalid substitution of a different family after taking the Fourier
limit. We state the join explicitly.

Fix a simple gauge group, a clip $p$ and path $R$, a bed $B$, and a
sequence of torus sides $n_i$ tending to infinity. Responses, their
fees, their three marked layers and greyness are as defined in
Section~\ref{sec:ex-annulus-fleet-faces}. We say that a response has the
\emph{tickets} of an analytic function $\mathcal U$ of the link field on
a torus if its zero-, one- and two-mark layers record exactly, on that
torus, the marked derivatives of $\mathcal U$: the derivatives at the
origin of $(y,z)\mapsto\mathcal U(e^ye^z)$ of order two in $y$ and of
order zero, one and two in $z$. Let $W$ be one fee-bounded analytic
Wilson response which is grey: it has the tickets of the Wilson action
at \emph{every} torus side, not only at the sides $n_i$. This is the
form in which the face estimates use it. Let $Q$ be one
fee-bounded correction response, with soft orders three, two, and one in
its zero-, one-, and two-mark layers, which at every chosen side $n_i$
has the tickets of the finite action of $p,R$ minus the Wilson action.
Assume, for this same bed and path: the path bounds of $R$ in $p$; the
attachment of $B$ to $R$; a metric strip bound for $B$; the floor of
$B,p$ with a constant $\epsilon>0$ at every torus side; and the sworn
condition at every chosen side $n_i$. Finally, the local-record family
of Section~\ref{sec:ex-annulus-fleet-faces} is formed with $V=W$ as its
vacuum response.
Assume $W$ and $Q$ have a common positive response strip width $\eta$ and
finite fee fixed before the odd side $L$ is chosen. This is the width used
for the uniform radial constants; it is distinct from the possibly smaller
strip $z_L$ of the regulated recording kernel.
For the logarithmic estimate, fix $0<\mu\leq1$ before applying the radial
bound, with $\mu=L^{-1}$ for the requested odd side $L$. The Wilson fee,
correction fee, floor parameter, and metric bound are common to the family;
the bed, clip, and side sequence are the same in the recording and face
estimates. The strip width and bound for the recorded kernel may depend on
this fixed $L$.

\begin{lemma}[Same-response compatibility]
\label{lem:ex-annulus-same-data}
Under these assumptions set $F=Q+W$. Then the following conclusions concern
this one $F$ and the same selected data:

\begin{enumerate}
\item For every chosen torus side, the zero-, one-, and two-mark tickets
  of $F$ are exactly the corresponding finite-action tickets. These are
  the three marked response layers used to identify the action variations
  in the finite test. If the regulated rank and
  reference matrices have their positive floors, the finite pole test is
  their logarithmic-determinant difference, and its conserving mixed
  coefficient is the site count times the normalized mesh mean of the
  shifted-fiber kernel of $F$.
\item The analytic kernel obtained from those shifted-fiber coefficients
  is independent of the torus side. Define the analytic picture by
  \[
    \mathcal Y_{ij}(p)=\frac12\int_{\mathbb T^4}F_{ij}(p,q)\,dq,
  \]
  with normalized Haar measure. At a conserving pair the exact
  finite-volume mixed derivative of the pole test is
  $2n_i^4\mathcal Y_{n_i,ij}(p)$, where
  \[
   \mathcal Y_{n,ij}(p)=\frac1{2n^4}
     \sum_{q\in(2\pi/n)(\mathbb Z/n\mathbb Z)^4}F_{ij}(p,q)
  \]
  is formed with the mesh mean; it is not $2n_i^4\mathcal Y_{ij}(p)$,
  the quantity with which $\operatorname{Err}$ compares it. The
  difference between the mesh picture and $\mathcal Y$ is controlled as
  follows. If the common kernel has bound $K$ on
  a strip of width $z>0$, the picture $\mathcal Y$ satisfies the
  recorded estimate
  \[
    \operatorname{Err}_{n_i}(\mathcal Y,\mathcal P_{n_i};
       \mathfrak r(z,K)e^{-zn_i/4})
  \]
  for every chosen side $n_i$.
\item For each color $c$, the radial bounds obtained from the same
  local-frame family and the same Wilson reference give
  \[
    |\alpha_c(\mathcal Y)+\lambda_G\log\mu^{-1}|\leq D_0.
  \]
  Here $D_0$ is fixed before the odd side parameter is chosen, while the
  strip constants $z,K$ and hence the recording error may depend on $L$.
\end{enumerate}
\end{lemma}
\begin{proof}
The fee bounds on $Q$ and $W$ restrict to their common smaller strip and add.
The exact response identities add as well: at each chosen side the zero-,
one-, and two-mark ticket of $Q$ is the finite-action ticket minus the
Wilson ticket, while the corresponding ticket of $W$ is the Wilson ticket.
Hence these three tickets of $F$ are exactly the finite-action tickets.
The order-zero remainder construction fixes the normalization; no extra
factor of the number of sites is inserted here.

For each side, the positive floors make the two determinant matrices
invertible at zero. Differentiating their logarithmic determinants gives
one direct second variation and the single ordered product of first
variations for each specified ordered pair. The ticket identity therefore
identifies that finite mixed coefficient with the conserving mesh sum of
the common shifted-fiber kernel. Nonconserving entries vanish by momentum
conservation. The fiber trace contributes a site count exactly once.
That the kernel is one function of the momenta, independent of the
side, is immediate, since the same $F$ is used. That it has a strip
bound is not: the kernel contains the inverses of the two regulated
denominators. The strip width $z$ and the bound $K$ are supplied by
Lemma~\ref{lem:ex-annulus-common-strip} below.
The full-mesh Fourier estimate of Lemma~\ref{lem:ex-pole-recording} now
applies to these exact finite identities, giving the displayed error at
every side, without a subsequence extraction.

For the third conclusion, apply the all-centered and shell face estimates
to the same selected local-frame family. The interacting and regulated
reference paths are the ones formed from $F$ and the fixed Wilson response;
the shell comparison uses the same bed and only momenta with positive
unregulated scale. The resulting two radial estimates are uniform in the
color and were fixed before choosing the odd side. Integrating over the
inner ball, the annulus, and the part outside the unit ball gives the
logarithmic estimate by Lemma~\ref{lem:path-radial-log}. The analytic
coefficient in that integral is $\alpha_c(\mathcal Y)$ for the picture
defined above: both it and the finite recorded coefficients are obtained
from the same shifted-fiber kernel, with the factor $1/2$ canceled by the
two-leg mixed derivative. This proves that the finite recording and the
radial drift belong to one response.
\end{proof}

\subsection{The common strip of the regulated kernel}

Fix $L$, put $\mu=L^{-1}\in(0,1]$, and let $F=Q+W$ be as above, with
all three layers holomorphic and periodic on a strip of width
$w>0$. Let the metric strip bound of $B$ hold at width $\sigma>0$ and
put
\[
 s=\min\{\sigma,w,\mu/128\}.
\]
In the notation of Section~\ref{sec:ex-annulus-fleet-faces}, the vector
total film is the film of $F$ plus $\mathsf B\mathsf B^{\rm r}$ plus
$R_{{\rm vec},\mu}$, and the ghost total film is
$Q_{\rm gh}+R_{{\rm gh},\mu}$. For a film with tab $a$ and slots
$b,c$, the two-leg figure is
\[
 \operatorname{fig}_{\xi\zeta}(p,q)=\operatorname{tr}\Bigl[
   a(q)^{-1}c_{\xi\zeta}(-p,p,q)
   -a(q)^{-1}b_\xi(-p,q+p)\,a(q+p)^{-1}b_\zeta(p,q)\Bigr],
\]
and the response kernel is one half of the figure of the vector total
film minus the figure of the ghost total film. This is the kernel
$F_{\xi\zeta}(p,q)$ of Lemma~\ref{lem:ex-pole-recording}: by the
determinant formula (RY8a), the conserving mixed coefficient of the pole
test on the torus of side $n_i$ is the site count times the grid mean
over $q$ of this kernel, once the tickets of the two total films are
known. The tickets of the vector total film are those of $F$, proved in
the first conclusion of Lemma~\ref{lem:ex-annulus-same-data}, together
with the tickets of the fixed films $\mathsf B$, $\mathsf D_a$ and of
their products, rears and inverses. The latter are an input of this
section, which we call \textup{(T)}: at every torus side, the slots of
$\mathsf B$ are the Fourier symbols of the covariant difference and of
its background derivatives, and the tickets of the products, rears and
inverses of the fixed films are those obtained by differentiating
products, transposes and inverses of matrices. The first assertion is a
direct evaluation on Fourier waves and the second a finite calculation,
but this ticket calculus is not written out in this text, and \textup{(T)}
is not established here.

The one fact about the selected path that enters is the value of the
zero-mark layer. By Lemma~\ref{lem:ex-soft-quadratic-hessian}, under the
path bounds, attachment, the floor at the side $n$ and the sworn
condition at the side $n$, the Hessian at zero of the quadratic part of
the finite action of $p,R$ on the torus of side $n$ is the kinetic form
of the bed,
$\sum_f\langle(d_1x)_f,(M_nd_1y)_f\rangle$ with $M_n$ the metric of $B$
on that torus. The cubic and quartic atom terms of the action vanish to
third order at zero and do not contribute to the Hessian. In Fourier
variables: at every grid momentum $k$ of the
torus of side $n_i$,
\begin{equation}\label{eq:ex-annulus-zero-layer}
 a_F(k)=\mathsf T\bigl[P_1(-k)^{\mathsf T}M(k)P_1(k)\bigr],
\end{equation}
where $a_F$ is the tab of the film of $F$, $M(k)$ the metric row of $B$
and $\mathsf T$, $P_1$ are as in (F14).

\begin{lemma}[Common strip]
\label{lem:ex-annulus-common-strip}
Under the standing assumptions of this section, with $L$ fixed, there
are $z>0$ and $K\geq0$ such that every entry of the response kernel is
holomorphic and periodic in $(p,q)$ on the strip of width $z$ and
bounded by $K$ there. The numbers $z,K$ do not depend on the torus
side; they may depend on $L$, on the bed and on the path.
\end{lemma}
\begin{proof}
\emph{Step 1: the denominator matching (F13) on the strip of width
$s$.} Write $\widehat P_j=\mathsf T(P_j)$ for the color-diagonal
extensions of the exterior symbols, which act on colored forms. The
tab of the vector total film is
$a_F(k)+\widehat P_0(k)\widehat P_0(-k)^{\mathsf T}+s_\mu(k)I$. Indeed
tabs multiply
under products of films and invert under inversion, so the tab of
$R_{{\rm gh},\mu}$ is $s_\mu I$ on colors, and the directional
sandwiches sum to $s_\mu I$. This function is holomorphic and periodic
on the strip of width $s$: $a_F$ by the fees, and $s_\mu$ because
$\omega(k)+\mu^2$ does not vanish on the strip of width $\mu/128$.
Indeed $\omega(k)=\sum_\nu(2-2\cos k_\nu)$, and for
$k_\nu=u+i\theta$ with $|\theta|\leq\mu/128\leq1$ one has
$\operatorname{Re}(2-2\cos k_\nu)=2-2\cos u\cosh\theta
\geq-2(\cosh\theta-1)\geq-2\theta^2$, so
$\operatorname{Re}(\omega(k)+\mu^2)\geq\mu^2-8(\mu/128)^2>0$. The
bed denominator $D^{\rm bed}(k)$ displayed in
Section~\ref{sec:ex-annulus-fleet-faces} is holomorphic and periodic on
the same strip, by the metric strip bound. By
\eqref{eq:ex-annulus-zero-layer} the two functions agree at every grid
momentum of every side $n_i$. Since $n_i\to\infty$, these grids are
dense in the real momentum torus, and both functions are continuous;
they agree at all real momenta. A holomorphic function on the strip
which vanishes at all real points vanishes identically, so they agree on
the strip. This is (F13).

\emph{Step 2: invertibility at real momenta.} Let $k$ be real and $x$ a
vector. The floor at every torus side gives, after passing from the
grids to all real momenta by density and continuity,
$\operatorname{Re}\langle v,M(k)v\rangle\geq\epsilon|v|^2$ for the
metric row. With the
exterior identity
$|\widehat P_1(k)x|^2+|\widehat P_0(k)^*x|^2=\omega(k)|x|^2$ for a
colored one-form $x$, and $s_\mu(k)\geq0$,
\[
 \operatorname{Re}\langle x,D^{\rm bed}(k)x\rangle
 \geq\epsilon|\widehat P_1(k)x|^2+|\widehat P_0(k)^*x|^2
      +s_\mu(k)|x|^2
 \geq\min(\epsilon,1)\bigl(\omega(k)+s_\mu(k)\bigr)|x|^2 .
\]
As shown in Section~\ref{sec:ex-annulus-fleet-faces},
$\omega+s_\mu\geq2^{-10}(\omega+\mu^2)\geq2^{-10}\mu^2>0$. Hence the
tab of the vector total film is invertible at every real momentum. The
tab of the ghost total film is the scalar $\omega+s_\mu$ on colors and
is invertible there as well.

\emph{Step 3: a complex neighborhood.} Let $a$ be either tab. The set
of points of the strip of width $s$ at which $\det a\ne0$ is open and
contains the compact real cell $[-\pi,\pi]^4$. It therefore contains
all $k$ with $\operatorname{Re}k\in[-\pi,\pi]^4$ and
$|\operatorname{Im}k_j|\leq r$ for some $0<r\leq s$, and by periodicity
the whole closed strip of width $r$. Take $r$ valid for both tabs.

\emph{Step 4: the kernel.} On the strip of width $r$ the inverse tabs
are holomorphic and periodic. The figure is a polynomial in the
entries of $a^{-1}$, $b$ and $c$ evaluated at $q$, $q+p$, $(-p,q+p)$,
$(p,q)$ and $(-p,p,q)$. If $(p,q)$ lies in the strip of width $r/2$,
all these arguments and their partial sums lie in the strip of width
$r$, where the three layers of both total films are holomorphic and
periodic. Hence every entry of the response kernel is holomorphic and
periodic on the joint strip of width $r/2$. Put $z=r/4$. The closed
strip of width $z$ is, modulo periods, the compact set
$[-\pi,\pi]^8\times[-z,z]^8$, on which the continuous entries are
bounded. Let $K$ be the sum of these bounds.
\end{proof}

Nothing in this proof is uniform in $L$: the lower bound $2^{-10}\mu^2$
degenerates as $\mu\to0$, and the radius $r$ comes from a compactness
argument. This is why the recording strip is chosen after $L$, while the
radial constants, which never use this strip, are chosen before it.

\subsection{The annulus estimate}

We can now state the output of
Sections~\ref{sec:ex-annulus-fleet-faces}--\ref{sec:ex-annulus-same-data}
in the form consumed by the one-step construction of
Section~\ref{subsec:path-stationary-chart}, where it is denoted
$\mathrm{HorizCalc}(G,\sigma,\epsilon,C,Z,\tau,A)$. It rests on two
inputs which are not proved in these sections.
\begin{itemize}
\item[(W)] \emph{Wilson response.} There are $k>0$, $E\geq0$ and a grey
 response $W$ with fees $(k,E)$.
\item[(S)] \emph{Correction service.} For all $\sigma,\tau>0$ and
 $C,Z,A\geq0$ and every real $\epsilon$ there are $v>0$ and
 $M\geq0$ such that for every sequence of sides $n_i\to\infty$, every
 bed $B$, clip $p$, path $R$ and number $\Lambda$ satisfying the
 hypotheses \eqref{eq:ex-annulus-hypotheses} below, there is a
 response $Q$ with fees $(v,M)$, soft of orders three, two and one in
 its three layers, which at every side $n_i$ has the tickets of the
 finite action of $p,R$ minus the Wilson action.
\end{itemize}
In the source, (W) is obtained from the exact Wilson coefficient
pictures of Proposition~\ref{prop:ex-wilson-stream}, and (S) from the
Wilson-subtracted pictures of Proposition~\ref{prop:ex-greyson}, whose
hypotheses are exactly \eqref{eq:ex-annulus-hypotheses}. In both cases
one converts exactly recorded normalized Taylor coefficient pictures of
orders two, three and four, for functions with vanishing first
derivative, into a response with fees whose three marked layers have
the corresponding tickets; in the case of (S) the pictures have ray
orders three, two and one and the layers are soft of the same orders,
while in the case of (W) no vanishing is required. The marked
derivatives are those of the split field
$\log(\exp x_+\exp x_-)$ used for the response arrays of
Section~\ref{sec:ex-soft-curvature-receipts}. That section states the
conversion for response arrays; its form for coefficient pictures with
fees, which is what (W) and (S) assert, is not stated as a result
there or elsewhere in this text. We therefore keep (W) and (S) as
explicit inputs.

\begin{proposition}[The annulus estimate]
\label{prop:ex-annulus-estimate}
Let $G$ be simple, let $\sigma,\epsilon,\tau>0$ and $C,Z,A\geq0$, and
assume \textup{(W)}, \textup{(S)} and the fixed-film ticket input
\textup{(T)} of this section. There is $D_0\geq0$ such that the
following holds for every odd integer $L\geq1$, every sequence of torus
sides $n_i\to\infty$, every bed $B$, clip $p$, path $R$ and number
$\Lambda$ such that
\begin{equation}\label{eq:ex-annulus-hypotheses}
\begin{gathered}
 R.\mathrm{Bounds}(p,\tau,A,\Lambda),\qquad
 \mathrm{Attached}(B,R),\qquad
 B.\mathrm{metric.strip}(\sigma,C),\qquad p.\mathrm{direct}\leq Z,\\
 \mathrm{Floor}(B,p,n,\epsilon,C)\ \hbox{for every torus side }n,
 \qquad \mathrm{Sworn}(B,n_i,G)\ \hbox{for every }i .
\end{gathered}
\end{equation}
There are a two-leg picture $\mathcal Y$, numbers $z>0$, $K\geq0$ and a
sequence $e_i\to0$ such that
\begin{enumerate}
\item every component of $\mathcal Y$ is holomorphic and periodic on the
 strip of width $z/2$ and bounded by $K$ there;
\item for every $i$, $\mathcal Y$ records the pole test
 $\mathcal P_{n_i}$ of $p,R$ with regulator $L^{-1}$ on the torus of
 side $n_i$, in the normalized sense \eqref{eq:output-record} with
 $n=2$, with error $n_i^4e_i$; here
 $e_i=\mathfrak r(z,K)e^{-zn_i/4}$;
\item $|\alpha_c(\mathcal Y)+\lambda_G\log L|\leq D_0$ for every color
 $c$.
\end{enumerate}
The constant $D_0$ depends only on $G,\sigma,\epsilon,C,Z,\tau,A$; the
picture $\mathcal Y$ and the numbers $z,K$ may depend on all the data,
including $L$. No limit in $i$ is part of the assertion.
\end{proposition}
\begin{proof}
Take $W,k,E$ from (W) and $v,M$ from (S); put $w=\min(k,v)$, so that
$W$ has fees $(w,E)$. Apply the conclusion (F11) of
Section~\ref{sec:ex-annulus-fleet-faces} with vacuum response $V=W$,
fee width $\eta=w$, floor constant $\epsilon$, bed bound $D=C$ and
correction fee $A=M$. It gives $a,b\geq0$, and we put
$D_0=c_{\rm wt}[a(\pi^2/2+|E_4|)+3\pi^2b]$. All of this precedes $L$
and the path data.

Now let $L$ and the data \eqref{eq:ex-annulus-hypotheses} be given, let
$Q$ be the response of (S), and put $F=Q+W$, $\mu=L^{-1}$ and
$s=\min\{\sigma,w,\mu/128\}$. The first conclusion of
Lemma~\ref{lem:ex-annulus-same-data} gives the tickets of $F$, and
Lemma~\ref{lem:ex-annulus-common-strip} gives $z,K$. Let $\mathcal Y$ be
the picture associated with the response kernel. The second conclusion
of Lemma~\ref{lem:ex-annulus-same-data}, that is
Lemma~\ref{lem:ex-pole-recording} in its normalized form
\textup{(RY4$'$)}, gives the first two assertions; $e_i\to0$ because
$z>0$ and $n_i\to\infty$.

For the third, observe that $(s,\mu,c,B,p,Q)$ is a local record in the
sense of Section~\ref{sec:ex-annulus-fleet-faces}, for every color $c$:
$Q$ has fees $(w,M)$ and the required soft orders; $0<\mu\leq1$ and
$0<s\leq\mu/128$; the layers of $F=Q+W$ are holomorphic on the strip of
width $s\leq w$ and the metric bound holds at width $s\leq\sigma$; the
denominator matching on the strip of width $s$ is Step~1 of the
preceding proof; and the floors hold at every side. This is the same
response $F$ and the same strip $s$ as in the recording. Hence (F11)
holds for its response coefficient with the constants $a,b$ chosen
above. The coefficient identity (AT1) applies to $\mathcal Y$ because
the kernel has the strip bound of
Lemma~\ref{lem:ex-annulus-common-strip}, and the radial integration in
the proof of Theorem~\ref{thm:ex-all-period-recording} gives the third
assertion with $\log\mu^{-1}=\log L$.
\end{proof}

Proposition~\ref{prop:ex-annulus-estimate} is
Theorem~\ref{thm:ex-all-period-recording} with its hypotheses on the
finite-layer realization, the common strip and the two-wave bounds
discharged; what remains are the inputs (W), (S) and (T) and, through
Lemma~\ref{lem:ex-soft-quadratic-hessian}, the quadratic coefficient of
the path action.

\paragraph{Selected construction of the common data.}
In the selected construction, the baseline $W$ is the fee-bounded Wilson
response. The correction $Q$ is the response produced from the attached
path and bed. Its finite tickets are the difference between the intended
action and the Wilson action. Combining them before evaluating any side
gives $F=Q+W$. The fee, ticket, floor, and metric conditions used above
are all on this combined response and the original selected bed.

The local-frame bounds are produced from that same Wilson reference: the
finite-volume vector and ghost responses, together with their one-leg and
two-leg extensions, are obtained from one family of local frames. The
all-centered estimate is applied at scale $|x|+\mu$; the shell estimate is
applied to the restriction $\mu\leq|x|$, where the unregulated reference
inverse exists. The two applications give the upper face and the lower
face-minus-halo price, respectively. These are precisely the hypotheses
used in the preceding face calculation; no change of bed, reference, or
finite-volume sequence is made between the face estimates and the ticket
identity.

In source terms, this is the composition of the Wilson baseline, the
path correction, ticket addition, finite determinant recording, and
uniform radial evaluation. The declaration-level correspondence and the
exact shared-object audit are in the companion evidence record.
